\section{Implementation}

GlassVM's shared crates define requests, preparation, event and evidence
types, normalizer contracts, the reference recorder, and the Python facade.
Each machine bundle implements the machine-specific descriptor, admission
rules, backend, session, and evidence adapters. Provider discovery in Python
uses installed entry points; the facade does not embed machine IDs or
emulator implementations.

The evaluation denominator contains three machine bundles. CHIP-8 is the
baseline reference bytecode virtual machine. PICO-8 and TIC-80 are independent
compatibility implementations with bounded Lua runtimes and documented
subsets. They do not claim exact upstream semantics. Table~\ref{tab:machines}
is generated from the checked-in qualitative machine manifest so that this
description does not duplicate machine-specific configuration details.

\begin{table}[t]
  \centering
  \small
  \caption{Qualitative machine heterogeneity and implementation boundaries.}
  \label{tab:machines}
  % Generated by generate_heterogeneity.py; edit heterogeneity.json instead.
\begin{tabular}{p{0.14\linewidth} p{0.28\linewidth} p{0.48\linewidth}}
\toprule
\textbf{Machine} & \textbf{Implementation role} & \textbf{Execution and fidelity boundary} \\
\midrule
CHIP-8 & Baseline reference bytecode virtual machine. & Native bytecode execution with machine-specific variants and configuration. \\
PICO-8 & Independent compatibility implementation. & Bounded Lua runtime implementing a documented headless subset; not a claim of exact upstream semantics. \\
TIC-80 & Independent compatibility implementation. & Bounded Lua runtime implementing a documented subset; non-Lua runtimes and unsupported upstream behavior remain outside its fidelity claim. \\
\bottomrule
\end{tabular}

\end{table}

The machine implementations share orchestration, preparation, evidence
delivery, and receipt contracts while retaining their own artifact, timing,
input, and runtime semantics. The evaluation uses that common boundary to
compare experimental procedures; it does not treat the three emulators as
semantically interchangeable machines.
